\section{Taxonomy}
\label{sec:taxonomy}

The taxonomy has three dimensions and one tag (Fig.~\ref{fig:taxonomy}), and it is organised around the
inferential claim a study makes. \textbf{Methodology} is primary and fixes the section in which a study
is discussed; \textbf{locus}, \textbf{concept provenance} and the \textbf{evidence type} tag are
recorded for every core study. The scheme maps the field's target space rather than the medical corpus
alone: some of its cells are populated mainly by general-domain work and hold few medical records, and
those cells are kept as gaps in the coverage map (Fig.~\ref{fig:corpus}b). The unit of coding is the
claim, one row with one evidence type. Table~\ref{tab:boundary} resolves the boundary cases.

\begin{figure*}[!tbp]
\centering
\includegraphics[width=0.96\linewidth]{fig4.png}
\caption{Taxonomy of feature-space interpretability methods for medical vision and multimodal foundation
models. The five columns are the method families of \S\ref{sec:methods}, assigned by the study's
inferential claim rather than its algorithm, with the sub-groups each section reviews. The band beneath
collects the intrinsic and concept-aligned designs, which are tagged rather than assigned a family. The
axis at the foot orders the families from representation structure and availability (predominantly
observational) to representation use and control (predominantly interventional); the evidence type of
any one claim is determined by that claim, not by its family.}
\label{fig:taxonomy}
\end{figure*}

\textbf{Methodology.} The five families are assigned by the \emph{primary inferential claim} of a
study, not by its algorithm, because one algorithm can support different claims. \emph{Structural and
geometric analysis} (\S\ref{sec:geometry}) characterises the shape of the space rather than a decoded
concept: image--text distance, spread, cross-model similarity, the change induced by adaptation, or the
structure of acquisition metadata. \emph{Supervised concept mapping and decodability}
(\S\ref{sec:decoding}) recovers an external concept from labelled data with a decoder, a direction or a
similarity structure. \emph{Feature decomposition and discovery} (\S\ref{sec:discovery}) breaks
representations into units, neurons, directions or dictionary features that are named afterwards.
\emph{Representation attribution, localisation and tracing} (\S\ref{sec:tracing}) assigns
representational content or output to a component, layer, token position or training example, without
intervening. \emph{Causal intervention and representation control} (\S\ref{sec:intervention}) changes
a representation and measures behaviour; the strength of the claim is set by the controls. Intrinsically
interpretable designs are properties of a model rather than a method: they take an
\emph{intrinsic-design} tag, and their claims stay in the analytic family, which is concept-layer
decodability in every case but one prototype design (\S\ref{sec:intrinsic}).

\begin{table*}[!tb]
\centering\small
\caption{Boundary cases and their assignment. \emph{Primary} determines the section in which a study is
discussed; \emph{secondary} is recorded in the coding.}
\label{tab:boundary}
\setlength\tabcolsep{4pt}\begin{tabular}{@{}L{5.4cm}L{2.8cm}L{1.9cm}L{5.2cm}@{}}
\toprule
Method as used in a study & Primary family & Secondary & Rule applied \\
\midrule
Principal-component, independent-component or non-negative matrix factorisation (PCA/ICA/NMF) of embeddings, directions then named & Discovery (\S\ref{sec:discovery}) & Geometry & Claim is about interpretable units, not about the shape of the space \\
PCA/effective rank/intrinsic dimension as a summary of the space & Geometry (\S\ref{sec:geometry}) & --- & No unit is interpreted \\
Representational similarity analysis against a concept-labelled stimulus set & Decoding (\S\ref{sec:decoding}) & Geometry & An external concept is decoded through the similarity structure \\
Cross-model CKA/RSA without concept labels & Geometry & --- & Structure only \\
Network dissection with a labelled concept dataset & Decoding & Discovery & Units are scored against supplied concepts \\
Sparse-autoencoder features that are then steered & Discovery (features) and Intervention (steering) & each other & Two claims, two sections, two evidence tags; each row records the other family as secondary \\
Concept erasure (iterative nullspace projection, INLP; least-squares concept erasure, LEACE) with a behavioural end-point & Intervention (\S\ref{sec:intervention}) & Decoding & Effect on behaviour is the claim \\
Concept erasure used only to test decodability & Decoding & --- & No behavioural end-point \\
Post-hoc concept bottleneck on a frozen encoder & Decoding & --- & Concept-layer decodability is the claim; the intrinsic-design tag records the model property (\S\ref{sec:intrinsic}) \\
Logit lens / tuned lens on decoder states & Tracing (\S\ref{sec:tracing}) & --- & Observational localisation of information across layers \\
Attention weights over image tokens & Tracing & --- & Component-level observation, not intervention (contextual when instrumental to a model or training-method paper, \S\ref{sec:method}) \\
Attention or token knockout & Intervention & Tracing & The representation is changed \\
Post-hoc operation on frozen embeddings (whitening, harmonisation, projection, transport, unlearning) with a downstream end-point & Intervention (inventory F9b) & Geometry or Decoding & The stored representation is changed and re-evaluated; a method that changes the encoder's training or weights only (fine-tuning, adversarial training) is contextual \\
Saliency map used to test a representation-level claim & Evaluation (\S\ref{sec:evaluation}) & --- & Readout-level instrument, not a feature-space method \\
\bottomrule
\end{tabular}
\end{table*}

\textbf{Locus.} A locus is the position along the computational path at which a representation is
taken: the vision encoder, the text encoder, the projected joint space of a dual encoder, the connector
output of a multimodal LLM, the decoder residual stream and its components, or the pre-readout
representation (\S\ref{sec:loci}, Fig.~\ref{fig:overview}). A study is coded by the locus at which it
measures or intervenes, not by where its claim is directed, and a study that repeats one measurement at
several loci is coded once, with the label ``Path''. The pre-readout representation, the connector
output and the text encoder carry almost no single-locus core claim (Fig.~\ref{fig:corpus}b); all six
loci are kept because that sparseness is itself a finding (\S\ref{sec:future}). Within the encoder the
coding also records the readout granularity as a \emph{sublocus}, which often decides what is decodable
(\S\ref{sec:findings}).

\textbf{Concept provenance.} Five sources are distinguished: \emph{expert- or ontology-defined}
vocabularies such as RadLex, SNOMED and the seven-point dermoscopy checklist; \emph{structured labels}
such as CheXpert and VinDr-CXR; \emph{corpus- or text-mined} concepts from reports or the literature;
\emph{model-generated} banks of class attributes; and \emph{none}, when units are discovered first and
labelled, if at all, afterwards.

\textbf{Evidence type.} A claim is \emph{observational} when it rests on measuring the representation
as the model computes it, and \emph{interventional} when it rests on a deliberate change, following the
first two levels of the causal hierarchy~\cite{bareinboim2022pearl}. The tag cuts across the families. It
separates the two statements most often run together in this literature, ``the model encodes $c$'' and
``the model uses $c$'', which rest on different
evidence~\cite{hewitt2019designing,elazar2020amnesic,ravfogel2020null}.

Fig.~\ref{fig:example} traces a hypothetical chest radiograph along the loci. It is a schematic of what
each family would report, not a summary of any one study, since only one study repeats a measurement at
successive loci of the same model~\cite{zhu2026lost}.

\begin{figure*}[!tbp]
\centering
\includegraphics[width=\linewidth]{fig5.png}
\caption{Hypothetical worked example. A chest radiograph with subtle blunting of the right costophrenic
angle passes through a frozen vision encoder and connector into visual tokens, which join the prompt in
the language model's residual stream and reach the pre-readout state and the generated report. Each
panel shows the artefact one family would return at its locus, the evidence type of that claim and the
controls it needs: (1) the geometry of the joint space; (2) a linear probe on the visual tokens
returning a concept score for right pleural effusion and a direction $v_c$; (3) a sparse dictionary
feature that fires on the costophrenic angle; (4) a map of which visual tokens and layers carry the
signal that reaches the report token ``effusion''; (5) the only interventional panel, adding
$\pm\alpha v_c$ to the residual stream and reading the change in the report against matched control
directions, sham edits, both signs and several magnitudes. No value is measured.}
\label{fig:example}
\end{figure*}
